\documentclass[11pt]{article}

\usepackage[margin=1.05in]{geometry}
\usepackage{amsmath,amssymb,amsthm}
\usepackage[T1]{fontenc}
\usepackage{lmodern}
\usepackage{microtype}
\usepackage{enumitem}
\usepackage{xcolor}

\setlength{\parindent}{0pt}
\setlength{\parskip}{0.55em}
\setlist[itemize]{leftmargin=1.5em,itemsep=0.2em,topsep=0.2em}

\newcommand{\A}{\mathcal A}
\newcommand{\B}{\mathcal B}
\newcommand{\sr}{\mathsf{single\_root}}
\newcommand{\st}{\mathsf{single\_term}}
\newcommand{\saa}{\mathsf{strong\_apder\_acc}}
\newcommand{\s}{\sigma_4}
\newcommand{\Sraw}{S}
\newcommand{\rsize}{\mathsf{rsize}}
\newcommand{\D}{D_1}
\newcommand{\RZERO}{\mathsf{0}}
\newcommand{\RONE}{\mathsf{1}}
\newcommand{\RCHAR}{\mathsf{char}}
\newcommand{\RSEQ}{\mathsf{seq}}
\newcommand{\RALTS}{\mathsf{alts}}
\newcommand{\RSTAR}{\mathsf{star}}

\begin{document}

{\Large\bfseries SEQ lane progress note}\\[-0.15em]
{\small Route-1 / \texttt{Posix\_Card\_Route1\_Seq}, private heap build green}

\vspace{0.6em}

\textbf{The invariant being paid is not a collapsed-continuation recurrence.}
The target carrier for the lane is the head contribution
\[
  \operatorname{card}\Big(
    \big(\sr(\RSEQ(h,t),k)\cup \st(h,\s(t,k))\big)
      - \big(\B k\cup \B(\s(t,k))\big)\Big)
  \le \rsize(h).                                      \tag{SEQ}
\]
The useful accounting charges the head size \(\rsize(h)\).  Charging
\(\D(h,\s(t,k))\) after the continuation has already collapsed is the false
recurrence: it loses the extra root row that appears in the \(\RSEQ\)-head
case.

\textbf{The RCHAR case is now a one-row proof.}
The term component is exactly the new boundary,
\[
  \st(\RCHAR(c),\s(t,k))=\B(\s(t,k)),
\]
so it disappears after subtracting
\(\B k\cup \B(\s(t,k))\).  What remains is the root row, and the checked
lemma is
\[
  \operatorname{card}\big(\sr(\RSEQ(\RCHAR(c),t),k)\big)\le 1.      \tag{RCHAR-root}
\]
Thus the branch closes immediately:
\begin{align*}
  &\operatorname{card}\Big(
    \big(\sr(\RSEQ(\RCHAR(c),t),k)\cup
          \st(\RCHAR(c),\s(t,k))\big)
      -\big(\B k\cup \B(\s(t,k))\big)\Big)\\
  &\hspace{35mm}\le \rsize(\RCHAR(c)).
\end{align*}
This was committed as the RCHAR root helper and the RCHAR head-core case.

\textbf{The star boundary is a one-step simulation, but only after subtraction.}
The tempting statement
\[
  \operatorname{card}\big(\B(\s(\RSTAR(r),k))\big)\le 1
\]
is false: when the star collapses to \(\RONE\), the whole old boundary
\(\B k\) can reappear.  The checked statement is the differential one:
\[
  \operatorname{card}\big(\B(\s(\RSTAR(r),k))-\B k\big)\le 1.       \tag{STAR-boundary}
\]
The proof first shows the row-level form
\[
  \operatorname{card}\Big(
      \operatorname{row}\big(\Sraw(\s(\RSTAR(r),k))\big)
      -\operatorname{row}\big(\Sraw(k)\big)\Big)\le 1,
\]
then uses the verified inclusion
\(\operatorname{row}(\Sraw(k))\subseteq \B k\).  The root identity also holds:
\[
  \sr(\RSTAR(r),k)=\B(\s(\RSTAR(r),k)).                            \tag{STAR-root}
\]

\textbf{The singleton-size ledger now has its star and sequence assembly bricks.}
For singleton carriers
\[
  \D(q,k)=\operatorname{card}\big(\A(\RALTS[q],k)-\B k\big),
\]
the atom leaves are green:
\[
  \D(\RZERO,k)\le \rsize(\RZERO),
  \qquad
  \D(\RCHAR(c),k)\le \rsize(\RCHAR(c)).
\]
The star recurrence is now checked:
\[
  \D(\RSTAR(r),k)
  \le
  1+\D\big(r,\s(\RSTAR(r),k)\big).                                \tag{D1-STAR}
\]
Hence any child bound
\(\D(r,\s(\RSTAR(r),k))\le \rsize(r)\)
immediately gives
\[
  \D(\RSTAR(r),k)\le \rsize(\RSTAR(r)).
\]
For a sequence head, the set algebra has also been isolated.  Let
\[
  c=\s(r_2,k).
\]
If the SEQ-head core for \(r_1\) is bounded by \(\rsize(r_1)\), and if the BND
lane supplies the boundary absorption
\[
  \operatorname{card}\big((\B c-\B k)\cup(\st(r_2,k)-\B k)\big)
  \le \D(r_2,k),
\]
then the checked assembly lemma gives
\[
  \D(\RSEQ(r_1,r_2),k)\le \rsize(r_1)+\D(r_2,k).                  \tag{D1-SEQ}
\]

\textbf{What remains is exactly the root-slack and cover work.}
The unclosed part is not another atom calculation.  It is the structural
root row in the \(\RSEQ\)-head case, plus the alternation branch cover:
\[
  \begin{aligned}
  \s(\RSEQ(r_1,r_2),k)&=\s(r_1,\s(r_2,k)),\\
  \sr(\RSEQ(\RSEQ(r_1,r_2),t),k)
  &\not\equiv
  \sr(\RSEQ(r_1,\RSEQ(r_2,t)),k)
  \end{aligned}
\]
definitionally.  That extra row is the genuine \(+1\) paid by
\(\rsize(\RSEQ(r_1,r_2))=1+\rsize(r_1)+\rsize(r_2)\), not a row to be erased
by reassociation.

The remaining imported dependencies should be stated explicitly when the main
lemma is assembled:
\[
  \text{BND: }\mathsf{boundary\_term\_absorb},
  \qquad
  \text{COVER: L1 singleton cover}.
\]
The current file is green under the private build command from
\texttt{ROUTE\_SEQ.md}; no shared heap wrapper was used.

\vfill
{\footnotesize
Last green commits: \texttt{0511e87} (RSEQ D1 assembly),
\texttt{ff7c8b6} (D1 star corollary),
\texttt{39e4fba} (D1 star step),
\texttt{62ca737} (star boundary shift),
\texttt{fbface5}/\texttt{6eff89b} (RCHAR core/root).
}

\end{document}
